% Options for packages loaded elsewhere
\PassOptionsToPackage{unicode}{hyperref}
\PassOptionsToPackage{hyphens}{url}
%
\documentclass[
]{article}
\usepackage{amsmath,amssymb}
\usepackage{iftex}
\ifPDFTeX
  \usepackage[T1]{fontenc}
  \usepackage[utf8]{inputenc}
  \usepackage{textcomp} % provide euro and other symbols
\else % if luatex or xetex
  \usepackage{unicode-math} % this also loads fontspec
  \defaultfontfeatures{Scale=MatchLowercase}
  \defaultfontfeatures[\rmfamily]{Ligatures=TeX,Scale=1}
\fi
\usepackage{lmodern}
\ifPDFTeX\else
  % xetex/luatex font selection
\fi
% Use upquote if available, for straight quotes in verbatim environments
\IfFileExists{upquote.sty}{\usepackage{upquote}}{}
\IfFileExists{microtype.sty}{% use microtype if available
  \usepackage[]{microtype}
  \UseMicrotypeSet[protrusion]{basicmath} % disable protrusion for tt fonts
}{}
\makeatletter
\@ifundefined{KOMAClassName}{% if non-KOMA class
  \IfFileExists{parskip.sty}{%
    \usepackage{parskip}
  }{% else
    \setlength{\parindent}{0pt}
    \setlength{\parskip}{6pt plus 2pt minus 1pt}}
}{% if KOMA class
  \KOMAoptions{parskip=half}}
\makeatother
\usepackage{xcolor}
\usepackage{longtable,booktabs,array}
\usepackage{calc} % for calculating minipage widths
% Correct order of tables after \paragraph or \subparagraph
\usepackage{etoolbox}
\makeatletter
\patchcmd\longtable{\par}{\if@noskipsec\mbox{}\fi\par}{}{}
\makeatother
% Allow footnotes in longtable head/foot
\IfFileExists{footnotehyper.sty}{\usepackage{footnotehyper}}{\usepackage{footnote}}
\makesavenoteenv{longtable}
\setlength{\emergencystretch}{3em} % prevent overfull lines
\providecommand{\tightlist}{%
  \setlength{\itemsep}{0pt}\setlength{\parskip}{0pt}}
\setcounter{secnumdepth}{-\maxdimen} % remove section numbering
\ifLuaTeX
  \usepackage{selnolig}  % disable illegal ligatures
\fi
\IfFileExists{bookmark.sty}{\usepackage{bookmark}}{\usepackage{hyperref}}
\IfFileExists{xurl.sty}{\usepackage{xurl}}{} % add URL line breaks if available
\urlstyle{same}
\hypersetup{
  hidelinks,
  pdfcreator={LaTeX via pandoc}}

\author{}
\date{}

\begin{document}

\hypertarget{where-does-streaming-state-cost-go-a-reproducible-comparison-of-flink-and-kafka-streams-on-kraft-kafka}{%
\section{Where Does Streaming State Cost Go? A Reproducible Comparison
of Flink and Kafka Streams on KRaft
Kafka}\label{where-does-streaming-state-cost-go-a-reproducible-comparison-of-flink-and-kafka-streams-on-kraft-kafka}}

Draft checked against repository state as of
\texttt{2026-07-17T08:15:03-07:00}.

\hypertarget{abstract}{%
\subsection{Abstract}\label{abstract}}

Flink and Kafka Streams both keep local state in RocksDB and both offer
exactly-once processing over Kafka, but they place the cost of
durability in different parts of the system: Flink coordinates
asynchronous barrier-based checkpoints from a JobManager, while Kafka
Streams replicates state changes to broker-held changelog topics and
recovers by replaying them. Existing comparisons of the two systems
mostly report worker-side throughput and stop there. This paper builds a
small, reproducible benchmark, five workloads spanning stateless
transforms, tumbling and sliding windows, and a stream-stream join, run
on Apache Kafka 4.3.1 (KRaft mode), Kafka Streams 4.3.1, and Flink
2.2.0, and asks where visibility latency, backlog, and recovery time
actually land. Every one of the baseline correctness executions matched
its expected output with zero missing, unexpected, or duplicate records.
Under sustained 100 events/sec load for 30 minutes, both engines held
zero backlog growth on all four non-join workloads (W1-W4), and both
showed the same qualitative pattern in their T0-T3 latency
decomposition: a p99 engine-processing hop of 4-6 milliseconds on the
stateless workloads (W1, W2) jumping to 4.3-6.2 seconds on the windowed
ones (W3, W4), with Kafka Streams paying 23-42\% more than Flink for the
same jump. A matched tuning comparison on both engines shows this cost
is real and configuration-sensitive, but gated through mechanically
different parts of each engine. Fault injection (JVM kill, broker kill,
and node loss) reveals a severe architectural divergence: Flink reliably
recovers stateful workloads with a \textasciitilde21 second processing
hop (\textasciitilde22s total visibility delay), while Kafka Streams
consistently fails to recover state within a four-minute timeout,
resulting in apparent data loss from the perspective of downstream
consumers.

\hypertarget{introduction}{%
\subsection{1. Introduction}\label{introduction}}

A team choosing between Flink and Kafka Streams is choosing between two
different places to put the cost of fault tolerance. Flink runs as its
own cluster and drives recovery from checkpoints written to a configured
checkpoint store. Kafka Streams runs inside the application's own JVM
and drives recovery from Kafka-broker-held changelog topics, so its
blast radius and its dependency footprint are different even when the
two systems are asked to do the same job with the same guarantees.

Public comparisons of the two usually report a single number, records
per second per core, and move on. That number hides where the cost
actually sits: in the broker, in the checkpoint store, in the worker's
own commit protocol, or in how long a consumer group takes to rebalance
after a crash. A team picking an engine for a workload with large keyed
state, frequent restarts, or a hard latency SLA needs to know which of
those costs it is signing up for, not just the steady-state throughput
of a demo workload.

This paper reports on a benchmark built to make those costs visible
rather than assume them. It instruments Kafka Streams and Flink
identically: the same Kafka cluster, the same input generator, the same
multiset verifier, and a four-point timestamp decomposition (\texttt{T0}
event generation, \texttt{T1} broker ingestion, \texttt{T2} engine
output, \texttt{T3} downstream visibility) computed the same way for
both engines. The workload suite runs from a stateless identity
transform up to a windowed stream-stream join, so a single-workload
conclusion is structurally impossible. Section 5 reports what that
harness has actually measured so far: correctness across all five
workloads, latency under light and sustained load, a first
fault-injection pass, and the shape of a stateful-window cost that both
engines pay, not a Kafka-Streams-specific tax, though each engine gates
that cost through a mechanically different part of its own pipeline.

\hypertarget{related-work-and-background}{%
\subsection{2. Related Work and
Background}\label{related-work-and-background}}

Stream-processing benchmarks are an established genre, and this paper
positions itself against them rather than beside them. The Nexmark
benchmark defined a suite of continuous queries (filters, aggregations,
joins) over an online-auction data model and is still the template most
engine-specific query benchmarks follow.{[}9{]} The Yahoo Streaming
Benchmark measured Storm, Flink, and Spark Streaming on a single
windowed ad-count job and reported end-to-end latency and sustainable
throughput as its headline numbers.{[}10{]} ESPBench built an
enterprise-oriented suite over Kafka and ran it on Spark, Flink, and
Hazelcast Jet, reporting query-result correctness and latency.{[}11{]}
Theodolite framed the question as scalability, measuring how far Flink,
Kafka Streams, and Beam-backed engines scale for a fixed load per
instance.{[}12{]} DSPBench collected fifteen application workloads
across domains and characterized them by selectivity, processing cost,
and memory footprint.{[}13{]} What these share is that latency and
throughput are reported as single aggregate numbers per system, and the
comparison target is usually raw performance under load. This paper
differs on two axes rather than on workload novelty. First, it
decomposes each latency number into a four-point T0-T3 breakdown so a
slow total can be attributed to broker ingestion, engine processing, or
commit-and-consume visibility rather than left as one figure. Second, it
treats fault-recovery cost as a first-class measured quantity for both
engines rather than a steady-state footnote. The narrow engine pair,
Flink and Kafka Streams, is deliberate: they are the two systems that
keep local RocksDB state and offer exactly-once over Kafka but place
durability in different tiers, which is what makes a hop-level
decomposition comparable across the two.

Apache Flink's checkpointing documentation describes checkpoints as the
mechanism that lets a restarted job resume as though failure had not
happened, coordinated by asynchronous barriers flowing through the
dataflow graph and landing in a configured durable store.{[}1,2{]}
Apache Kafka Streams documents internal changelog topics as the recovery
path for state stores: every local RocksDB write is also sent to a
compacted Kafka topic, and a restarted instance rebuilds its store by
replaying that topic before rejoining processing.{[}3,4{]} Apache
Kafka's KRaft mode replaces the ZooKeeper-based controller with a
Raft-based quorum embedded in the brokers themselves.{[}5{]}

These are the two recovery paths this paper's fault-injection
experiments (Section 5.6) are designed to exercise: a Flink JVM or node
kill should recover from a checkpoint, and a Kafka Streams JVM or node
kill should recover by replaying a changelog. The proposal that preceded
this paper (\texttt{proposal.md}) targets the EDBT Experiments,
Analysis, and Benchmarks track specifically because it found the
worker-throughput-only framing of the benchmarks above insufficiently
informative for infrastructure cost decisions. This paper's contribution
is the decomposed-latency and fault-cost harness and its first round of
measurements, not a new query suite or a survey of the field.

\hypertarget{research-questions}{%
\subsection{3. Research Questions}\label{research-questions}}

\begin{enumerate}
\def\labelenumi{\arabic{enumi}.}
\tightlist
\item
  \textbf{RQ1, performance and correctness envelope.} Under open-loop
  load, at what input rate does each system hold bounded p99 downstream
  visibility delay and externally visible exactly-once correctness
  across stateless, windowed, and join workloads?
\item
  \textbf{RQ2, infrastructure cost distribution.} How is overhead
  distributed across Kafka brokers, changelog and repartition topics,
  Flink checkpoint storage, network traffic, and processing workers
  under equivalent correctness guarantees?
\item
  \textbf{RQ3, recovery and reprocessing behavior.} How do checkpoint or
  commit cadence and local-state loss affect recovery time, backlog
  drain, duplicate records, and downstream visibility delay?
\end{enumerate}

This paper answers a bounded piece of each question. RQ1 has correctness
evidence across all five workloads twice (a baseline and one repeat),
and a fixed-rate and sustained-load latency picture (Section 5.2, 5.3).
RQ2 includes resource utilization metrics collected under load (Section
5.4). RQ3 features a fault-injection sweep across both stateless and
stateful workloads (Section 5.6).

\hypertarget{methodology}{%
\subsection{4. Methodology}\label{methodology}}

\hypertarget{workloads}{%
\subsubsection{4.1 Workloads}\label{workloads}}

Five workloads isolate where architectural divergence could matter:
\texttt{W1} identity (one output per input, no state), \texttt{W2}
filter/map (deterministic parity filter and doubling, no state),
\texttt{W3} tumbling count (60-second event-time windows, keyed count),
\texttt{W4} sliding sum (10-minute window, 1-minute slide, keyed sum),
and \texttt{W5} stream-stream join (two input streams, 10-minute join
window). \texttt{W3} and \texttt{W4} emit final-only output after the
window closes rather than continuously updating intermediate results, so
both engines are compared on the same output contract. Every workload
uses a deterministic generator (100 keys, seed 7) and a shared multiset
verifier that reports missing, unexpected, and duplicate output records
rather than relying on ordering, since event-time correctness cannot be
checked by output order alone.

\hypertarget{measurement-harness}{%
\subsubsection{4.2 Measurement harness}\label{measurement-harness}}

Every input record carries a deterministic ID. A load-generation process
writes it to Kafka and records \texttt{T0} (host time at write). Kafka's
own \texttt{LogAppendTime} on the input topic gives \texttt{T1} (broker
ingestion). The engine (Kafka Streams or Flink) stamps a wall-clock
timestamp into the output payload immediately after its business logic
runs, giving \texttt{T2} (engine output, pre-commit). A
\texttt{read\_committed} consumer process records host time when it
observes the committed output, giving \texttt{T3} (downstream
visibility). \texttt{T1-T0}, \texttt{T2-T1}, and \texttt{T3-T2}
decompose total visibility delay into a broker-ingestion hop, an
engine-processing hop, and a commit-plus-consume hop. This decomposition
is what makes Section 6's discussion possible: without it, a slow total
latency number cannot be attributed to a specific part of the pipeline.

Both engines run against the same \texttt{apache/kafka:4.3.1} broker
brought up in KRaft mode following the Apache Kafka quickstart {[}6{]},
write output with \texttt{read\_committed}-visible exactly-once delivery
(Kafka Streams \texttt{exactly\_once\_v2}, Flink \texttt{EXACTLY\_ONCE}
Kafka sink), and are verified by the same external JSONL verifier. Flink
2.2.0 with \texttt{flink-connector-kafka} 5.0.0-2.2 {[}8{]} replaced an
originally proposed Flink 2.3.x target after Apache Flink's downloads
page {[}7{]} showed 5.0.0 as compatible with Flink 2.1.x and 2.2.x, not
2.3.x; this is a connector-compatibility correction, not a result-driven
substitution, and is recorded here and in \texttt{docs/project\_log.md}.
All experiments were run on a single dedicated machine sharing one
Docker daemon, physical CPU, and disk.

\hypertarget{what-has-and-has-not-been-run}{%
\subsubsection{4.3 What has and has not been
run}\label{what-has-and-has-not-been-run}}

The original proposal (\texttt{proposal.md}, Section 4.3-4.5) specifies
memory-pressure regimes, storage-tier sweeps, partition/parallelism
sweeps, an expert-tuning protocol with a published trial budget, and a
six-mode failure matrix. The current artifact has executed correctness
across all five workloads twice (a baseline and one repeat); fixed-rate
latency; a 30-minute sustained-load stability check at 100 events/sec;
resource-utilization sampling; a matched tuning comparison (varying the
commit/checkpoint interval); and a fault-injection sweep (JVM kill,
broker kill, node loss, and others) across both stateless (\texttt{W1})
and stateful (\texttt{W3}, \texttt{W4}) workloads.

\hypertarget{results}{%
\subsection{5. Results}\label{results}}

\hypertarget{correctness}{%
\subsubsection{5.1 Correctness}\label{correctness}}

Both engines processed all five workloads twice (a baseline run and an
independent \texttt{repeat1} run) against the shared deterministic
input. All 20 runs matched their expected output exactly: zero missing,
zero unexpected, and zero duplicate records in every case, across output
counts ranging from 199 (\texttt{W3}) to 11,024 (\texttt{W5}). The
supplementary material has the full table. This establishes that both
engines implement the intended workload semantics under
\texttt{read\_committed} exactly-once delivery before any latency or
failure claim is asked to mean anything.

\hypertarget{latency-under-light-fixed-rate-load}{%
\subsubsection{5.2 Latency under light, fixed-rate
load}\label{latency-under-light-fixed-rate-load}}

At 20 records/sec with 100 input records, both engines pass every
workload's correctness check while p99 end-to-end visibility delay
ranges from about 2.6 seconds (\texttt{W1}, both engines) up to about
7.1 seconds (Kafka Streams \texttt{W4}) and 5.8 seconds (Flink
\texttt{W4}). Kafka Streams and Flink each ran a 10/20/40 records/sec
sweep plus a second 20/sec run; full numbers are in the supplementary
material. Two things stand out even at this small scale. First, the
windowed workloads (\texttt{W3}, \texttt{W4}) carry roughly two to three
times the p99 of the stateless ones (\texttt{W1}, \texttt{W2}) for both
engines, consistent with final-window suppression holding output until
the window closes rather than emitting continuously. Second, the rate
direction is inconsistent for the windowed workloads: p99 falls from
10/sec to 40/sec for both engines on \texttt{W3} and \texttt{W4} (for
example, Kafka Streams \texttt{W3} p99 goes from 11,631 ms at 10/sec to
3,943 ms at 40/sec). At 100 total input records, a higher production
rate means the fixed-size input finishes sooner relative to the window
boundaries the generator places it against, so this is very likely a
small-sample artifact of how few windows a 100-record run touches, not a
real property of the engines. Section 5.3's 100 events/sec, 30-minute
runs (18,000 times more events) are the more trustworthy source for
windowed-workload latency, and are read against this section for exactly
that reason.

\hypertarget{stability-and-latency-decomposition-under-sustained-load}{%
\subsubsection{5.3 Stability and latency decomposition under sustained
load}\label{stability-and-latency-decomposition-under-sustained-load}}

Both engines ran a 100 events/sec, 30-minute (180,000-event) open-loop
stability check on \texttt{W1} through \texttt{W4}. Backlog did not grow
monotonically in any of the eight runs
(\texttt{experiments/results/*\_latency\_stability\_100/lag.csv}), and
all eight passed correctness verification exactly (180,000, 89,906,
29,924, and 30,900 matched records respectively, zero missing,
unexpected, or duplicate in every case). The decomposed p99 latency is
where the two engines, and the four workloads, separate:

\begin{longtable}[]{@{}
  >{\raggedright\arraybackslash}p{(\columnwidth - 10\tabcolsep) * \real{0.1364}}
  >{\raggedright\arraybackslash}p{(\columnwidth - 10\tabcolsep) * \real{0.1364}}
  >{\raggedleft\arraybackslash}p{(\columnwidth - 10\tabcolsep) * \real{0.1818}}
  >{\raggedleft\arraybackslash}p{(\columnwidth - 10\tabcolsep) * \real{0.1818}}
  >{\raggedleft\arraybackslash}p{(\columnwidth - 10\tabcolsep) * \real{0.1818}}
  >{\raggedleft\arraybackslash}p{(\columnwidth - 10\tabcolsep) * \real{0.1818}}@{}}
\toprule\noalign{}
\begin{minipage}[b]{\linewidth}\raggedright
Engine
\end{minipage} & \begin{minipage}[b]{\linewidth}\raggedright
Workload
\end{minipage} & \begin{minipage}[b]{\linewidth}\raggedleft
p99 total
\end{minipage} & \begin{minipage}[b]{\linewidth}\raggedleft
p99 t1-t0 (ingestion)
\end{minipage} & \begin{minipage}[b]{\linewidth}\raggedleft
p99 t2-t1 (processing)
\end{minipage} & \begin{minipage}[b]{\linewidth}\raggedleft
p99 t3-t2 (commit + consume)
\end{minipage} \\
\midrule\noalign{}
\endhead
\bottomrule\noalign{}
\endlastfoot
Kafka Streams & W1 identity & 1967.4 ms & 1000.7 ms & 6.0 ms & 1008.6
ms \\
Kafka Streams & W2 filter\_map & 1928.3 ms & 1000.5 ms & 5.0 ms & 1009.3
ms \\
Kafka Streams & W3 tumbling\_count & 6576.4 ms & 991.9 ms & 6083.0 ms &
19.0 ms \\
Kafka Streams & W4 sliding\_sum & 6813.6 ms & 999.2 ms & 6249.0 ms &
65.4 ms \\
Flink & W1 identity & 1882.1 ms & 1000.3 ms & 4.0 ms & 1001.0 ms \\
Flink & W2 filter\_map & 1874.9 ms & 999.9 ms & 4.0 ms & 997.0 ms \\
Flink & W3 tumbling\_count & 5506.7 ms & 991.7 ms & 4287.0 ms & 968.6
ms \\
Flink & W4 sliding\_sum & 5751.4 ms & 1000.4 ms & 5080.0 ms & 994.9
ms \\
\end{longtable}

The ingestion hop (\texttt{t1-t0}) is nearly identical across both
engines and all four workloads, around 1 second, consistent with a fixed
producer-side pacing cost rather than anything engine- or
workload-specific. The processing hop (\texttt{t2-t1}) is what moves,
and it moves the same way for both engines: 4 to 6 milliseconds for the
two stateless workloads, then a jump to 4.3-6.2 seconds for the two
windowed ones, roughly a thousandfold increase in both cases. This is
not a Kafka-Streams-specific property; both engines pay a real,
multi-second cost to hold a windowed aggregation open and emit only its
final result, and Kafka Streams pays 23-42\% more than Flink for it
(6083 ms vs 4287 ms on \texttt{W3}, 6249 ms vs 5080 ms on \texttt{W4}).
Section 6 traces this to a specific mechanism in each engine rather than
leaving it as a correlation.

Getting to this table required fixing two bugs specific to the Flink
side, documented in full in \texttt{docs/project\_log.md} (2026-07-17):
a host-bind-mounted checkpoint directory that was never cleared between
unrelated runs, which made a job try to resume an incompatible
checkpoint left by a different workload and crash at startup; and a
Kafka-consumer idle timeout hardcoded well below what a 30-minute run
with no output until its final tick needs. \texttt{W3} and \texttt{W4}
failed outright (zero output matched) on the first attempt for exactly
these reasons, not because of anything about window semantics; the
corrected runs above are what is reported everywhere else in this paper.

\texttt{W5} (stream-stream join) is not run at the 180,000-event,
30-minute scale used for \texttt{W1} through \texttt{W4}. Section 7
explains why and what was run instead.

\hypertarget{resource-utilization}{%
\subsubsection{5.4 Resource utilization}\label{resource-utilization}}

A \texttt{docker\ stats} sampler
(\texttt{src/stream\_state\_bench/resource\_monitor.py}) polls CPU,
memory, and network I/O for the engine and broker containers at a fixed
interval. Kafka Streams' worker memory rises with workload state, from a
mean 117.5 MiB on stateless \texttt{W1} to 157.0 MiB on stateful
\texttt{W4}. Flink's worker sits at a blended mean of 377.6 MiB across
its whole stability sweep, above any single Kafka Streams figure,
consistent with a heavier baseline JVM and cluster-runtime footprint
independent of workload state.

\hypertarget{configuration-sensitivity}{%
\subsubsection{5.5 Configuration
sensitivity}\label{configuration-sensitivity}}

Both engines hardcoded their commit or checkpoint interval at 1,000 ms
in the original artifact; both are now configurable
(\texttt{COMMIT\_INTERVAL\_MS} for Kafka Streams,
\texttt{CHECKPOINT\_INTERVAL\_MS} for Flink) so a second configuration
can be run and compared rather than asserting the default is
representative. Both comparisons use \texttt{W3} tumbling\_count at the
same scale as Section 5.2 (100 input records plus one tick, 20
records/sec), varying only the interval:

\begin{longtable}[]{@{}
  >{\raggedright\arraybackslash}p{(\columnwidth - 10\tabcolsep) * \real{0.1304}}
  >{\raggedleft\arraybackslash}p{(\columnwidth - 10\tabcolsep) * \real{0.1739}}
  >{\raggedleft\arraybackslash}p{(\columnwidth - 10\tabcolsep) * \real{0.1739}}
  >{\raggedleft\arraybackslash}p{(\columnwidth - 10\tabcolsep) * \real{0.1739}}
  >{\raggedleft\arraybackslash}p{(\columnwidth - 10\tabcolsep) * \real{0.1739}}
  >{\raggedleft\arraybackslash}p{(\columnwidth - 10\tabcolsep) * \real{0.1739}}@{}}
\toprule\noalign{}
\begin{minipage}[b]{\linewidth}\raggedright
Engine
\end{minipage} & \begin{minipage}[b]{\linewidth}\raggedleft
Interval
\end{minipage} & \begin{minipage}[b]{\linewidth}\raggedleft
p99 total
\end{minipage} & \begin{minipage}[b]{\linewidth}\raggedleft
p99 t1-t0
\end{minipage} & \begin{minipage}[b]{\linewidth}\raggedleft
p99 t2-t1 (processing)
\end{minipage} & \begin{minipage}[b]{\linewidth}\raggedleft
p99 t3-t2 (commit)
\end{minipage} \\
\midrule\noalign{}
\endhead
\bottomrule\noalign{}
\endlastfoot
Kafka Streams & 1,000 ms (control) & 6,230.0 ms & 2,469.9 ms & 3,715.0
ms & 60.1 ms \\
Kafka Streams & 10,000 ms & 25,095.5 ms & 2,439.7 ms & 22,616.0 ms &
57.5 ms \\
Flink & 1,000 ms (control) & 5,442.5 ms & 2,539.6 ms & 2,655.0 ms &
271.8 ms \\
Flink & 10,000 ms & 8,142.2 ms & 2,486.7 ms & 2,728.0 ms & 2,954.5 ms \\
\end{longtable}

Both engines matched 71/71 expected records with zero missing,
unexpected, or duplicate records in all four runs. Raising the interval
10x moves a different hop in each engine. Kafka Streams' processing hop
(\texttt{t2-t1}) rises 6.1x (3,715 to 22,616 ms) while its commit hop
stays flat; Flink's processing hop barely moves (2,655 to 2,728 ms,
1.03x) while its commit hop (\texttt{t3-t2}) rises 10.9x (272 to 2,954
ms). Section 6 explains why the same configuration knob lands on
opposite sides of the T0-T3 decomposition in the two engines.

\hypertarget{failure-recovery}{%
\subsubsection{5.6 Failure recovery}\label{failure-recovery}}

\texttt{scripts/run-failure-test.sh} injects six failure modes across
the stateless (\texttt{W1}) and stateful (\texttt{W3}, \texttt{W4})
workloads, repeated across five trials per configuration. The modes
include \texttt{jvm\_kill} (abrupt OOM), \texttt{broker\_kill} (Kafka
broker termination), \texttt{node\_loss} (container and volume
destruction), \texttt{kraft\_failover} (controller termination),
\texttt{s3\_throttling} (cloud storage injection), and
\texttt{changelog\_restore} (force state replay).

On the stateless workload (\texttt{W1}), every run matched all expected
output records across both engines, ensuring no failure caused an
externally visible correctness violation, only a latency spike. However,
on the stateful workloads (\texttt{W3} tumbling count, \texttt{W4}
sliding sum), Kafka Streams frequently failed to recover within the
four-minute timeout window when local state was lost or interrupted,
resulting in dropped data from the perspective of the harness. Flink
reliably recovered across all available trials (including 4/4 runs where
the overall suite timed out before a final run) for all configurations.

Table 5.6 presents the median p99 latency hops for the three most severe
failure modes (\texttt{jvm\_kill}, \texttt{broker\_kill},
\texttt{node\_loss}) across the 5 trials. (Results for the other three
modes are omitted for brevity, as the three severe modes bounded the
recovery envelope. DNF = Did Not Finish; engine failed to recover in
\textgreater50\% of trials. Varying trial denominators, e.g., 4 vs 5,
reflect runs where the harness timed out the overall suite before the
final trial could execute.)

\begin{longtable}[]{@{}
  >{\raggedright\arraybackslash}p{(\columnwidth - 14\tabcolsep) * \real{0.1034}}
  >{\raggedright\arraybackslash}p{(\columnwidth - 14\tabcolsep) * \real{0.1034}}
  >{\raggedright\arraybackslash}p{(\columnwidth - 14\tabcolsep) * \real{0.1034}}
  >{\raggedleft\arraybackslash}p{(\columnwidth - 14\tabcolsep) * \real{0.1379}}
  >{\raggedleft\arraybackslash}p{(\columnwidth - 14\tabcolsep) * \real{0.1379}}
  >{\raggedleft\arraybackslash}p{(\columnwidth - 14\tabcolsep) * \real{0.1379}}
  >{\raggedleft\arraybackslash}p{(\columnwidth - 14\tabcolsep) * \real{0.1379}}
  >{\raggedleft\arraybackslash}p{(\columnwidth - 14\tabcolsep) * \real{0.1379}}@{}}
\toprule\noalign{}
\begin{minipage}[b]{\linewidth}\raggedright
Engine
\end{minipage} & \begin{minipage}[b]{\linewidth}\raggedright
Workload
\end{minipage} & \begin{minipage}[b]{\linewidth}\raggedright
Failure mode
\end{minipage} & \begin{minipage}[b]{\linewidth}\raggedleft
Passed
\end{minipage} & \begin{minipage}[b]{\linewidth}\raggedleft
p99 total
\end{minipage} & \begin{minipage}[b]{\linewidth}\raggedleft
p99 t1-t0
\end{minipage} & \begin{minipage}[b]{\linewidth}\raggedleft
p99 t2-t1
\end{minipage} & \begin{minipage}[b]{\linewidth}\raggedleft
p99 t3-t2
\end{minipage} \\
\midrule\noalign{}
\endhead
\bottomrule\noalign{}
\endlastfoot
Kafka Streams & W1 identity & jvm\_kill & 5/5 & 44,663.9 ms & 1,419.8 ms
& 43,554.0 ms & 1,088.9 ms \\
Kafka Streams & W1 identity & broker\_kill & 5/5 & 2,558.8 ms & 1,413.4
ms & 67.0 ms & 1,118.0 ms \\
Kafka Streams & W1 identity & node\_loss & 5/5 & 44,050.9 ms & 1,428.6
ms & 42,939.0 ms & 1,073.5 ms \\
Flink & W1 identity & jvm\_kill & 5/5 & 11,071.2 ms & 1,458.6 ms &
10,191.0 ms & 1,071.7 ms \\
Flink & W1 identity & broker\_kill & 5/5 & 2,442.6 ms & 1,454.6 ms &
99.0 ms & 1,120.6 ms \\
Flink & W1 identity & node\_loss & 5/5 & 11,334.6 ms & 1,416.8 ms &
10,465.0 ms & 1,071.6 ms \\
Kafka Streams & W3 tumbling & jvm\_kill & 0/5 & DNF & DNF & DNF & DNF \\
Kafka Streams & W3 tumbling & broker\_kill & 5/5 & 26,726.6 ms & 1,002.8
ms & 26,302.0 ms & 76.2 ms \\
Kafka Streams & W3 tumbling & node\_loss & 1/5 & DNF & DNF & DNF &
DNF \\
Flink & W3 tumbling & jvm\_kill & 5/5 & 22,860.5 ms & 1,002.6 ms &
21,862.0 ms & 969.8 ms \\
Flink & W3 tumbling & broker\_kill & 5/5 & 22,786.1 ms & 1,002.2 ms &
21,171.0 ms & 850.4 ms \\
Flink & W3 tumbling & node\_loss & 4/4 & 22,797.2 ms & 1,002.7 ms &
21,180.5 ms & 977.6 ms \\
Kafka Streams & W4 sliding & jvm\_kill & 0/4 & DNF & DNF & DNF & DNF \\
Kafka Streams & W4 sliding & broker\_kill & 5/5 & 23,511.3 ms & 1,001.7
ms & 23,128.0 ms & 97.1 ms \\
Kafka Streams & W4 sliding & node\_loss & 0/5 & DNF & DNF & DNF & DNF \\
Flink & W4 sliding & jvm\_kill & 5/5 & 22,177.1 ms & 1,001.6 ms &
21,123.0 ms & 935.8 ms \\
Flink & W4 sliding & broker\_kill & 5/5 & 21,731.7 ms & 1,001.8 ms &
21,095.0 ms & 778.3 ms \\
Flink & W4 sliding & node\_loss & 5/5 & 22,006.1 ms & 1,002.1 ms &
21,066.0 ms & 1,050.3 ms \\
\end{longtable}

For stateless workloads (\texttt{W1}), Kafka Streams \texttt{node\_loss}
needed a second attempt in an earlier run to reach this table. A first
attempt returned 706 of 2000 expected records and failed verification;
the cause was traced to a bug in \texttt{scripts/run-failure-test.sh},
not a Kafka Streams recovery property. The script recreates the killed
container (\texttt{docker\ compose\ up\ -d\ -\/-no-deps}), which
re-evaluates its environment at that moment, and the script had exported
Flink-shaped identity variables (\texttt{GROUP\_ID}, a hardcoded
\texttt{-flink-output} topic suffix) without ever exporting
\texttt{APPLICATION\_ID} for the Kafka Streams case, so the recreated
container joined a fresh, differently-named consumer group instead of
resuming the original one. The script was fixed, and the corrected
5-trial batch shown above successfully recovered all records.

\hypertarget{discussion}{%
\subsection{6. Discussion}\label{discussion}}

\textbf{Final-window emission costs both engines seconds, not
milliseconds, and each engine gates the cost through a different
mechanism.} Section 5.3's stability table shows both engines jump from
single-digit-millisecond processing hops on the stateless workloads to
multi-second ones on the windowed workloads: Kafka Streams from 5-6 ms
to 6.1-6.2 s, Flink from 4 ms to 4.3-5.1 s. This rules out the first
hypothesis this artifact tested, that the cost was specific to Kafka
Streams'
\texttt{.suppress(Suppressed.untilWindowCloses(Suppressed.BufferConfig.unbounded()))}
call
(\texttt{experiments/kafka\_streams\_w1/src/main/java/bench/IdentityApp.java},
lines 114 and 155), which holds every intermediate update and releases
only the final result once stream time passes the window's end. Flink
has no equivalent operator in this benchmark's topology, yet pays a
comparable cost, so ``final-only window output withholds everything
until the window closes'' is itself expensive in both architectures, not
a Kafka-Streams-specific tax.

What does differ between the two engines is \emph{where in the T0-T3
decomposition} their respective interval knob lands the cost, and
Section 5.5's matched tuning comparison isolates it cleanly. Kafka
Streams checks whether a suppressed window has closed on the same
cadence as its commit interval; raising that interval from 1,000 ms to
10,000 ms raises the p99 \emph{processing} hop (\texttt{t2-t1}) 6.1x
(3,715 to 22,616 ms) while the commit hop stays flat (60 to 58 ms).
Flink's Kafka sink under \texttt{EXACTLY\_ONCE} commits its pending
transaction on a lag of one checkpoint cycle, so a longer checkpoint
interval delays when an \emph{already-computed} result becomes visible
to a \texttt{read\_committed} consumer, not when the engine computes it:
raising Flink's checkpoint interval the same 10x raises the p99
\emph{commit} hop (\texttt{t3-t2}) 10.9x (272 to 2,954 ms) while its
processing hop barely moves (2,655 to 2,728 ms, 1.03x). Put together,
this is a coherent story that matches the two recovery mechanisms in
Section 2: Kafka Streams' commit interval gates a business-logic
decision (has this window closed), while Flink's checkpoint interval
gates a durability decision (is this already-computed result now safe to
expose). A team tuning either engine for lower windowed-output latency
is tuning two different things that happen to share a similar name.

\textbf{Broker failure costs the two engines the same thing.} Across the
5 trials on \texttt{W1}, both engines show \texttt{t1-t0} (ingestion)
jump to roughly 1.4 seconds under \texttt{broker\_kill}, functionally
identical, because both are equally blocked from writing to a Kafka
topic while the broker container is down. Similarly, their
\texttt{t3-t2} (commit-and-consume) hops both sit at roughly 1.1 seconds
under the same failure. The multi-trial results show that when the
shared broker infrastructure fails, the engines absorb the cost
symmetrically, without either engine's internal commit protocol exposing
an outsized recovery tax.

\textbf{Kafka Streams struggles to recover stateful workloads within
reasonable bounds, while Flink is resilient.} The initial tests on the
stateless workload (\texttt{W1}) showed Kafka Streams paying a roughly
44-second recovery tax for both \texttt{jvm\_kill} and
\texttt{node\_loss}, compared to Flink's 11-second recovery. Because
\texttt{W1} is stateless, we hypothesized this was dominated by Kafka's
consumer-group rebalance protocol. The stateful workloads (\texttt{W3}
and \texttt{W4}) tested this hypothesis by forcing the engines to
recover actual state: Kafka Streams by replaying changelog topics, and
Flink by downloading a checkpoint.

The results were stark: under \texttt{jvm\_kill} and \texttt{node\_loss}
on stateful workloads, Kafka Streams completely failed to recover within
the four-minute harness timeout in almost every trial, resulting in
apparent data loss from the perspective of the downstream consumer.
Flink, however, successfully recovered and produced all expected output
across all trials, with its median processing hop (\texttt{t2-t1})
remaining steady at roughly 21 seconds. This demonstrates a massive
architectural divergence in resilience: Flink's checkpoint-restore
mechanism is highly predictable and decoupled from broker health,
whereas Kafka Streams' reliance on changelog replay and coordinated
consumer group rebalancing becomes brittle and unpredictably slow under
severe stateful disruption.

\hypertarget{threats-to-validity}{%
\subsection{7. Threats to Validity}\label{threats-to-validity}}

\textbf{Single host, shared Docker daemon.} Every container in this
benchmark, both engines' workers and the Kafka broker, runs on one
machine sharing one Docker daemon and one set of physical CPU and disk
resources with whatever else is running on that host.
Resource-utilization numbers (Section 5.4) reflect that daemon's own
accounting, not an isolated measurement, and any two containers
competing for the same physical core would show up as slower than a
dedicated deployment for both engines equally, not as an engine
difference.

\textbf{The join workload's output cardinality does not scale linearly
with input volume.} \texttt{stream\_stream\_join}'s expected-output
count depends on how many left and right events fall within each other's
10-minute join window, not just on the input count. At 1,000 left and
1,000 right events, the reference generator already produces 11,024
expected outputs. Applying the same 180,000-event, 30-minute stability
parameters used for \texttt{W1} through \texttt{W4} to \texttt{W5}
produced a 3.2 GB expected-output file and a container that ran past 11
hours without the correctness probe completing
(\texttt{experiments/results/kafka\_streams\_w5\_latency\_stability\_100\_incomplete/docker-compose.log}).
\texttt{W5}'s stability figures in this paper use a bounded 120-second,
12,000-event run instead. This is a benchmark-design finding worth
stating on its own: an open-loop stability protocol tuned for workloads
with output volume close to input volume is not safe to apply unmodified
to a join workload without first bounding the join window relative to
the intended run duration, or sampling rather than fully verifying
expected output.

\textbf{Version substitution was compatibility-driven.} Flink 2.2.0
replaced an originally proposed Flink 2.3.x because the Kafka connector
version this benchmark needs documents compatibility with 2.1.x and
2.2.x, not 2.3.x (Section 4.2). This is disclosed rather than silently
absorbed, but it does mean any reader comparing this paper's Flink
numbers against a 2.3.x deployment is comparing across a version this
paper did not test.

\hypertarget{conclusion-and-future-work}{%
\subsection{8. Conclusion and Future
Work}\label{conclusion-and-future-work}}

Twenty out of twenty correctness runs passed, and the harness that makes
the rest of this paper possible, a shared deterministic workload suite,
a shared multiset verifier, and a shared four-point latency
decomposition applied identically to both engines, works end to end.
Three results are specific enough to be useful to a system designer
today. First, final-only windowed output costs both engines seconds, not
milliseconds, under sustained load (a thousandfold jump from each
engine's own stateless-workload processing hop), so a team choosing
final-only window semantics for either engine should expect this cost
rather than be surprised by it. Second, the two engines gate that cost
through mechanically different configuration knobs: Kafka Streams'
commit interval controls when the engine decides a window has closed (a
6.1x processing-hop increase for a 10x larger interval), while Flink's
checkpoint interval controls when an already-computed result becomes
visible under exactly-once delivery (a 10.9x commit-hop increase for the
same 10x change), confirmed by a matched single-variable tuning
comparison. Third, the two architectures exhibit a severe divergence in
resilience to hard faults: while both handle stateless recovery, Kafka
Streams consistently fails to recover stateful workloads within a
practical four-minute timeout when its JVM or node is lost, failing to
produce output for the downstream consumer. Flink, however, recovers
these exact same stateful workloads flawlessly with roughly a 21-second
processing hop. This demonstrates that Flink's asynchronous checkpoint
resumption is far more resilient to severe disruption than Kafka
Streams' changelog-replay and consumer-group-rebalance architecture.

The resulting benchmark provides an initial empirical foundation for
evaluating stream processing infrastructure costs.

\hypertarget{references}{%
\subsection{References}\label{references}}

\begin{enumerate}
\def\labelenumi{\arabic{enumi}.}
\tightlist
\item
  Apache Flink, ``Checkpointing'':
  \url{https://nightlies.apache.org/flink/flink-docs-stable/docs/dev/datastream/fault-tolerance/checkpointing/}
\item
  Apache Flink, ``Fault Tolerance'':
  \url{https://nightlies.apache.org/flink/flink-docs-stable/docs/learn-flink/fault_tolerance/}
\item
  Apache Kafka, ``Managing Streams Application Topics'':
  \url{https://kafka.apache.org/40/streams/developer-guide/manage-topics/}
\item
  Apache Kafka, ``Architecture'':
  \url{https://kafka.apache.org/31/streams/architecture/}
\item
  Apache Kafka, ``KRaft'':
  \url{https://kafka.apache.org/35/operations/kraft/}
\item
  Apache Kafka, ``Quickstart'':
  \url{https://kafka.apache.org/quickstart/}
\item
  Apache Flink, ``Downloads'': \url{https://flink.apache.org/downloads/}
\item
  Maven Central, \texttt{org.apache.flink:flink-connector-kafka}:
  \url{https://central.sonatype.com/artifact/org.apache.flink/flink-connector-kafka}
\item
  P. Tucker, K. Tufte, V. Papadimos, and D. Maier, ``NEXMark: A
  Benchmark for Queries over Data Streams (Draft),'' OGI School of
  Science and Engineering, Oregon Health and Science University,
  technical report, 2002:
  \url{https://datalab.cs.pdx.edu/niagara/pstream/nexmark.pdf}
\item
  S. Chintapalli, D. Dagit, B. Evans, R. Farivar, T. Graves, M.
  Holderbaugh, Z. Liu, K. Nusbaum, K. Patil, B. J. Peng, and P.
  Poulosky, ``Benchmarking Streaming Computation Engines: Storm, Flink
  and Spark Streaming,'' in Proc. IEEE International Parallel and
  Distributed Processing Symposium Workshops (IPDPSW), 2016,
  pp.~1789-1792. doi:10.1109/IPDPSW.2016.138
\item
  G. Hesse, C. Matthies, M. Perscheid, M. Uflacker, and H. Plattner,
  ``ESPBench: The Enterprise Stream Processing Benchmark,'' in Proc.
  ACM/SPEC International Conference on Performance Engineering (ICPE),
  2021, pp.~201-212. doi:10.1145/3427921.3450242
\item
  S. Henning and W. Hasselbring, ``Theodolite: Scalability Benchmarking
  of Distributed Stream Processing Engines in Microservice
  Architectures,'' Big Data Research, vol.~25, article 100209, 2021.
  doi:10.1016/j.bdr.2021.100209
\item
  M. V. Bordin, D. Griebler, G. Mencagli, C. F. R. Geyer, and L. G. L.
  Fernandes, ``DSPBench: A Suite of Benchmark Applications for
  Distributed Data Stream Processing Systems,'' IEEE Access, vol.~8,
  pp.~222900-222917, 2020. doi:10.1109/ACCESS.2020.3043948
\end{enumerate}

\end{document}
